\section{问题一的建立与求解}

267名孕妇共1082条纵向观测数据中，每名孕妇平均包含约4个时间点的重复测量，个体内相关性对参数估计有直接影响，需要在建模中加以处理。

\subsection{混合模型框架的选择依据}

为评估个体层面聚类效应的强度，首先计算组内相关系数（ICC）。以$\sigma^2_b$表示个体间方差、$\sigma^2_\varepsilon$表示残差方差，ICC定义为
\begin{equation}
  \mathrm{ICC} = \frac{\sigma^2_b}{\sigma^2_b + \sigma^2_\varepsilon}
\end{equation}

计算结果为$\mathrm{ICC}=0.60$，Y染色体浓度总变异的60\%来源于个体间差异。若采用合并OLS回归忽略这一聚类结构，标准误将被系统性低估，假阳性膨胀。混合模型框架通过随机效应显式建模个体异质性，可保证纵向重复测量下统计推断的有效性。

\subsection{广义加性混合模型的建立}

探索性分析表明，Y染色体浓度随孕周的变化呈非线性趋势，线性假设拟合不足。本文建立\textbf{广义加性混合模型（GAMM，记为M1）}作为主模型，同时建立\textbf{线性混合效应模型（LME，记为M2）}作为基准，通过嵌套比较评估非线性成分的贡献。

GAMM模型（M1）的具体形式为
\begin{equation}\label{eq:gamm}
  y_{ij} = \beta_0 + s_1(t_{ij}) + s_2(\mathrm{BMI}_i) + b_{0i} + b_{1i}\,t_{ij} + \varepsilon_{ij}
\end{equation}

其中$y_{ij}$为第$i$名孕妇在第$j$次观测时的Y染色体浓度，$t_{ij}$为对应孕周。各成分含义如下。

\begin{itemize}
  \item $s_1(\cdot)$为孕周的非线性平滑函数，采用B样条基函数，自由度$\mathrm{df}=4$；
  \item $s_2(\cdot)$为BMI的非线性平滑函数，采用B样条基函数，自由度$\mathrm{df}=3$；
  \item $b_{0i} \sim N(0,\,\sigma^2_{b0})$为随机截距，刻画个体基线浓度差异；
  \item $b_{1i} \sim N(0,\,\sigma^2_{b1})$为孕周的随机斜率，刻画个体浓度增长速率差异；
  \item $\varepsilon_{ij} \sim N(0,\,\sigma^2_\varepsilon)$为观测误差。
\end{itemize}

LME基准模型（M2）将两个平滑函数替换为线性项
\begin{equation}\label{eq:lme}
  y_{ij} = \beta_0 + \beta_1\,t_{ij} + \beta_2\,\mathrm{BMI}_i + b_{0i} + b_{1i}\,t_{ij} + \varepsilon_{ij}
\end{equation}

M2与M1共享相同的随机效应结构，二者构成嵌套关系，可通过似然比检验直接比较。

\subsubsection{模型求解步骤}

模型估计按以下步骤执行。

\textbf{Step 1}（数据预处理）\quad 对267名孕妇的1082条观测记录进行缺失值检查与异常值筛查，确认数据完整性后按孕妇编号排列，以适配混合模型的分组结构。

\textbf{Step 2}（参数估计）\quad 采用限制最大似然法（REML）估计方差成分，固定效应参数通过L-BFGS优化器求解。REML在方差分量估计上具有无偏性，适合样本量有限的纵向数据。

\textbf{Step 3}（模型比较）\quad 对M1和M2分别以ML重新拟合（保证似然比检验的有效性），计算似然比统计量$\Lambda = 2(\ell_{\mathrm{M1}} - \ell_{\mathrm{M2}})$，在$\chi^2$分布下检验非线性成分的显著性。

\textbf{Step 4}（残差诊断）\quad 对最优模型的残差进行正态性检验（Shapiro-Wilk）和分布形态分析（偏度、峰度），验证模型假设的合理性。

\subsection{模型求解结果}

表~\ref{tab:q1-compare}汇总了M1（GAMM）与M2（LME）在主要评价指标上的对比结果。

\begin{table}[H]
  \centering
  \caption{GAMM（M1）与LME（M2）模型评价指标对比}\label{tab:q1-compare}
  \begin{tabular}{lcccc}
    \toprule
    模型 & AIC & BIC & RMSE & 对数似然 \\
    \midrule
    M1（GAMM） & $-5240.80$ & $-5180.96$ & 0.0101 & 2632.40 \\
    M2（LME）  & $-5174.04$ & $-5139.13$ & 0.0111 & --- \\
    \bottomrule
  \end{tabular}
\end{table}

似然比检验结果为$\chi^2=76.76$（$\mathrm{df}=5$，$p<10^{-15}$），M1在统计意义上优于M2。M1的AIC较M2低66.76，BIC低41.83，RMSE降低约9\%，三项指标一致支持非线性模型。

M1的固定效应估计如下。截距$\hat{\beta}_0=0.0107$；孕周B样条基系数依次为$0.0153$、$0.0173$、$0.0434$、$0.1238$，呈递增趋势，反映Y染色体浓度随孕周加速上升；BMI的B样条基系数为$0.1023$、$0.0127$、$0.0150$。随机效应方面，随机截距方差$\hat{\sigma}^2_{b0}=0.001886$，随机斜率方差$\hat{\sigma}^2_{b1}=0.000008$，个体基线差异远大于斜率差异，与ICC=0.60的结果一致。

残差诊断显示偏度为$0.125$，峰度为$4.15$，Shapiro-Wilk统计量$W=0.943$。残差分布接近对称，峰度略高于正态分布的理论值3，存在轻微尖峰，但在$n=1082$的样本量下对REML估计的稳健性影响有限。

图~\ref{fig:q1-partial}展示了M1与M2在孕周和BMI两个维度上的偏效应对比。GAMM的孕周偏效应曲线呈非线性上升形态，孕晚期阶段斜率增大明显，LME仅能以固定斜率的直线近似这一趋势。

\begin{figure}[H]
  \centering
  \includegraphics[width=0.85\textwidth]{Q1/experiments/round1/figures/M1_vs_M2_partial_effects.png}
  \caption{M1（GAMM）与M2（LME）偏效应对比}\label{fig:q1-partial}
\end{figure}

图~\ref{fig:q1-gest}给出GAMM模型中孕周偏效应的单独可视化及其95\%置信带。B样条基系数从早期的$0.0153$递增至$0.1238$，Y染色体浓度在孕早期增长平缓，中晚期后加速上升。这一非线性特征是GAMM相对于LME取得改进的核心来源。

\begin{figure}[H]
  \centering
  \includegraphics[width=0.75\textwidth]{Q1/experiments/round1/figures/M1_GAMM_partial_gest_week.png}
  \caption{GAMM模型孕周偏效应曲线}\label{fig:q1-gest}
\end{figure}

\subsection{结果分析}

GAMM的求解结果揭示了三条规律。

（1）孕周对Y染色体浓度具有强非线性正向效应。样条基系数序列$0.0153 \to 0.0173 \to 0.0434 \to 0.1238$表明浓度增速随孕周推进逐步加快，似然比检验（$\chi^2=76.76$，$p<10^{-15}$）证实这一非线性成分不可省略。

（2）BMI对Y染色体浓度存在负向调节作用。BMI样条基的首项系数$0.1023$后迅速回落至$0.0127$和$0.0150$，高BMI孕妇体内母体血浆体积增大，对胎儿游离DNA产生稀释效应，检测浓度偏低。

（3）个体异质性显著。ICC=0.60意味着超过半数的浓度变异归因于个体间差异，随机截距方差$\hat{\sigma}^2_{b0}=0.0019$远大于随机斜率方差$\hat{\sigma}^2_{b1}=0.000008$，不同孕妇的基线浓度水平差异是个体异质性的主要表现。

从NIPT临床实践角度看，孕周效应的非线性意味着对所有孕妇施行统一采样策略并非最优。孕早期浓度增速慢，过早采样可能因浓度不足而增加检测失败风险；孕中晚期浓度快速攀升，灵敏度相对充裕。BMI的稀释效应进一步加剧了这一问题，高BMI孕妇在相同孕周下浓度系统性偏低，需要推迟采样以补偿浓度劣势。这两项发现自然引出一个后续问题，即如何根据BMI水平制定分组化的最优检测时点，这正是问题二要回答的内容。
